\documentclass[12pt,a4paper]{article}
\usepackage[UTF8]{ctex}
\usepackage{amsmath,amssymb,amsfonts,geometry,enumitem}
\geometry{left=2.5cm,right=2.5cm,top=2.5cm,bottom=2.5cm}
\title{2024年新高考Ⅰ卷·数学 考场作答（考生B）}
\author{考生：李欣怡\quad 准考证号：202406071315}
\date{2024年6月7日}
\begin{document}
\maketitle
\begin{center}
{\large \textbf{（本答卷为中等偏上考生，部分题有失误，解答过程不完整）}}
\end{center}

\section*{一、选择题（1~8题）}
\begin{enumerate}[label=\arabic*.]
\item A \quad （解不等式，选A）
\item C \quad （复数运算，算得1-i）
\item D \quad （向量垂直，x=2）
\item A \quad （用两角和差，蒙的A，应该对吧）
\item B \quad （侧面积相等，算体积得$3\sqrt3\pi$）
\item C \quad （分段单调，没太懂，蒙了C）
\item C \quad （画图数交点，6个）
\item B \quad （递推推了几项，选B）
\end{enumerate}

\section*{二、多项选择题（9~11题）}
\begin{enumerate}[label=\arabic*.]
\item BC \quad （正态分布，A错B对，C对D错）
\item AC \quad （A对，B试了不对，C对，D没验证，就选AC）
\item BC \quad （看不懂图，按感觉选BC）
\end{enumerate}

\section*{三、填空题（12~14题）}
\begin{enumerate}[label=\arabic*.]
\item $\dfrac{3}{2}$ \quad （通径和勾股，算得$\frac32$）
\item $\ln 3$ \quad （切线斜率2，设切点算出来$\ln3$，不知道对不对，填了$\ln3$）
\item $\dfrac{7}{16}$ \quad （概率题，枚举法，得$\frac{7}{16}$）
\end{enumerate}

\section*{四、解答题}
\begin{enumerate}[label=\arabic*.]
\item （13分）
\begin{enumerate}[label=(\arabic*)]
\item 由余弦定理得$C=45^\circ$，$\sin C=\sqrt2\cos B\Rightarrow \cos B=\frac12$，$B=60^\circ$。
\item $A=75^\circ$，面积$S=\frac12 ac\sin B=3+\sqrt3$，得到$ac=4\sqrt3+4$，再用正弦定理，我算得$c=2\sqrt2$。（过程写了，应该对）
\end{enumerate}

\item （15分）
\begin{enumerate}[label=(\arabic*)]
\item $b^2=9$，$a^2=12$，$e=\frac12$。
\item 设直线$l$斜率$k$，用面积公式，我解得$k=\frac32$，另一个没算出来，就写了一个$y=\frac32x-3$。（估计扣一半分）
\end{enumerate}

\item （15分）
\begin{enumerate}[label=(\arabic*)]
\item 证明$AD\parallel$平面$PBC$，写了因为$AD\perp PB$且$PA\perp AD$，所以$AD\perp$平面$PAB$，得$AD\perp AB$，又$AD\subset$底面，$BC\subset$底面，由$AB\perp BC$？不太对，这题证得不好。
\item 建系后算$AD$，设$D(t,0,0)$，列方程解得$t=3$，所以$AD=3$。（过程有点乱，答案写3）
\end{enumerate}

\item （17分）
\begin{enumerate}[label=(\arabic*)]
\item $a$最小值为$-2$，求导做出来了。
\item 证明中心对称：$f(x)+f(2-x)=2a$，所以关于$(1,a)$对称，写了。
\item 求$b$范围，不会做，只写了个$b\ge -\frac23$，没有过程。
\end{enumerate}

\item （17分）
\begin{enumerate}[label=(\arabic*)]
\item $(i,j)=(1,2),(1,6),(5,6)$，枚举出来的。
\item 不会证明，写“由等差数列可构造”几个字。
\item 空白。
\end{enumerate}

\vspace{1cm}
\begin{flushright}
\textbf{估分：110~115分} \quad 时间不够，后面大题没写完。
\end{flushright}
\end{document}